\section{引言：对齐研究的时代背景}

本次讲座由 Nathan Lambert 主讲，他曾在 UC Berkeley 完成博士学位，先后在 Hugging Face 和 Allen Institute for AI（AI2）工作，是后训练（post-training）和对齐（alignment）领域的核心研究者之一。讲座标题"Life after DPO"试图回答一个关键问题：\textbf{在 DPO 之后，对齐研究的下一步方向是什么？}

\begin{figure}[H]
\centering
\includegraphics[width=0.95\textwidth]{slides-images/slide-000.jpg}
\caption{讲座封面：Life after DPO，Nathan Lambert，Allen Institute for AI\protect\footnotemark}
\end{figure}
\footnotetext{Slides 第1页。}

讲座的核心结构分为三部分：
\begin{enumerate}
    \item \textbf{历史回顾}：从语言模型发展到 RLHF 和 DPO 的演进
    \item \textbf{评估工具}：RewardBench——如何评估奖励模型
    \item \textbf{前沿探索}：DPO vs PPO 的实证比较，以及在线（online）方法的兴起
\end{enumerate}

\begin{knowledgebox}{讲者背景}
Nathan Lambert 的研究背景是强化学习（RL）。他从机器人 RL 转向语言模型 RL，在 Hugging Face 期间参与了 Zephyr 等开源对齐模型的开发，后加入 AI2 主导 T\"ulu 系列模型和 RewardBench 项目。他既有学术研究视角，也有工业界一线经验，对开放对齐生态有深刻理解。
\end{knowledgebox}

\subsection{语言模型的简要历史}

Lambert 快速回顾了语言模型的发展脉络：

\begin{itemize}
    \item \textbf{2017}：Transformer 论文发表
    \item \textbf{2018}：GPT-1、ELMo、BERT——奠定现代 NLP 基础
    \item \textbf{2019--2020}：GPT-2 与 Scaling Laws——人们开始意识到大规模语言模型的潜力
    \item \textbf{2021}：Stochastic Parrots 论文——在 ChatGPT 之前就对语言模型提出警告
    \item \textbf{2022年末}：ChatGPT 发布——本应是 OpenAI 的一个低调 demo，却引爆了全球关注
\end{itemize}

\begin{figure}[H]
\centering
\includegraphics[width=0.95\textwidth]{slides-images/slide-001.jpg}
\caption{语言模型发展的简要时间线\protect\footnotemark}
\end{figure}
\footnotetext{Slides 第2页。}

\begin{importantbox}{ChatGPT 能否脱离 RLHF 而存在？}
Lambert 提出了一个核心问题：ChatGPT 能否在没有 RLHF 的情况下存在？答案是：RLHF 似乎是\textbf{必要但不充分}的（necessary but not sufficient）。预训练提供了基础能力，但 RLHF 和后训练技术是让模型真正"可用"的关键一步。
\end{importantbox}

\begin{figure}[H]
\centering
\includegraphics[width=0.95\textwidth]{slides-images/slide-005.jpg}
\caption{RLHF 是必要但不充分的条件\protect\footnotemark}
\end{figure}
\footnotetext{Slides 第6页。}

Lambert 特别引用了 Meta 在 Llama 2 技术报告中的一段话："强化学习以其不稳定性著称，对 NLP 研究社区来说像是某种神秘领域；然而，强化学习被证明非常有效，尤其考虑到其成本和时间效率。"这段话生动地反映了 2023 年年中时业界对 RLHF 从怀疑到认可的态度转变。

\begin{figure}[H]
\centering
\includegraphics[width=0.95\textwidth]{slides-images/slide-006.jpg}
\caption{Anthropic Constitutional AI 论文中的迭代改进曲线：展示了多轮 RLHF 如何逐步提升模型性能\protect\footnotemark}
\end{figure}
\footnotetext{Slides 第7页。}

\subsection{本章小结}

语言模型的发展经历了从自回归损失函数到 Transformer 再到大规模预训练的漫长积累。ChatGPT 的成功不仅依赖于强大的预训练基础，更需要 RLHF 等后训练技术的加持。理解这一历史背景，是理解当前对齐研究方向的前提。

%% ========== Part 2: 核心概念与定义 ==========

